\documentclass[aspectratio=169, xcolor=dvipsnames]{beamer}

% --- Theme Setup ---
\usetheme{Madrid}
\usecolortheme{default}

% --- Packages ---
\usepackage{graphicx}
\usepackage{xcolor}
\usepackage{booktabs}
\usepackage{hyperref}
\usepackage{adjustbox}
\usepackage{amssymb}
\usepackage{amsmath}
\usepackage{listings}
\usepackage{caption}

\lstset{
  basicstyle=\ttfamily\small,
  keywordstyle=\color{blue}\bfseries,
  stringstyle=\color{red},
  showstringspaces=false,
  breaklines=true
}

% --- Title Information ---
\title{Module 12}
\subtitle{Intermediate SQL/1}
\author{Milav Dabgar}
\institute{www.milav.in}
\date{\today}

\begin{document}

% Title Slide
\begin{frame}
    \titlepage
\end{frame}

\begin{frame}{Module Recap}
    \begin{itemize}
        \item SQL Examples Practiced
    \end{itemize}
\end{frame}

\begin{frame}{Module Objectives}
    \begin{itemize}
        \item To understand nested subquery in SQL
        \item To understand processes for data modification
    \end{itemize}
\end{frame}

\begin{frame}{Module Outline}
    \begin{itemize}
        \item Nested Subqueries
        \item Modifications of the Database
    \end{itemize}
\end{frame}

\section{Nested Subqueries}

\begin{frame}
  \vfill
  \centering
  \begin{beamercolorbox}[sep=8pt,center,shadow=true,rounded=true]{title}
    \usebeamerfont{title}Nested Subqueries\par%
  \end{beamercolorbox}
  \vfill
\end{frame}

\begin{frame}[fragile]{Nested Subqueries}
    \begin{itemize}
        \item SQL provides a mechanism for the nesting of subqueries
        \item A subquery is a \texttt{select}-\texttt{from}-\texttt{where} expression that is nested within another query
        \item The nesting can be done in the following SQL query
        \begin{lstlisting}[language=SQL]
select A1, A2, ..., An 
from r1, r2, ..., rm
where P
        \end{lstlisting}
        \textbf{as follows:}
        \begin{itemize}
            \item \texttt{Ai} can be replaced by a subquery that generates a single value
            \item \texttt{ri} can be replaced by any valid subquery
            \item \texttt{P} can be replaced with an expression of the form:
            \item \texttt{B <operation> (subquery)}
            \item where \texttt{B} is an attribute and \texttt{<operation>} to be defined later
        \end{itemize}
    \end{itemize}
\end{frame}

\section{Subqueries in the Where Clause}

\begin{frame}
  \vfill
  \centering
  \begin{beamercolorbox}[sep=8pt,center,shadow=true,rounded=true]{title}
    \usebeamerfont{title}Subqueries in the Where Clause\par%
  \end{beamercolorbox}
  \vfill
\end{frame}

\begin{frame}{Subqueries in the Where Clause}
    \begin{itemize}
        \item Typical use of subqueries is to perform tests:
        \begin{itemize}
            \item For set membership
            \item For set comparisons
            \item For set cardinality
        \end{itemize}
    \end{itemize}
\end{frame}

\begin{frame}[fragile]{Set Membership}
    \begin{itemize}
        \item Find courses offered in Fall 2009 and in Spring 2010. (\texttt{intersect} example)
        \begin{lstlisting}[language=SQL, basicstyle=\ttfamily\scriptsize]
select distinct course_id 
from section 
where semester = 'Fall' and year = 2009 and 
      course_id in 
          (select course_id 
           from section 
           where semester = 'Spring' and year = 2010);
        \end{lstlisting}
        \item Find courses offered in Fall 2009 but not in Spring 2010. (\texttt{except} example)
        \begin{lstlisting}[language=SQL, basicstyle=\ttfamily\scriptsize]
select distinct course_id 
from section 
where semester = 'Fall' and year = 2009 and 
      course_id not in 
          (select course_id 
           from section 
           where semester = 'Spring' and year = 2010);
        \end{lstlisting}
    \end{itemize}
\end{frame}

\begin{frame}[fragile]{Set Membership (2)}
    \begin{itemize}
        \item Find the total number of (distinct) students who have taken course sections taught by the instructor with ID 10101
        \begin{lstlisting}[language=SQL, basicstyle=\ttfamily\scriptsize]
select count(distinct ID) 
from takes 
where (course_id, sec_id, semester, year) in 
          (select course_id, sec_id, semester, year 
           from teaches 
           where teaches.ID = 10101);
        \end{lstlisting}
        \item Note: Above query can be written in simpler manner. The formulation above is simply to illustrate SQL features.
    \end{itemize}
\end{frame}

\begin{frame}[fragile]{Set Comparison - 'some' Clause}
    \begin{itemize}
        \item Find names of instructors with salary greater than that of \textbf{some} (at least one) instructor in the Biology department
        \begin{lstlisting}[language=SQL, basicstyle=\ttfamily\scriptsize]
select distinct T.name 
from instructor as T, instructor as S 
where T.salary > S.salary and S.dept_name = 'Biology';
        \end{lstlisting}
        \item Same query using \texttt{some} clause
        \begin{lstlisting}[language=SQL, basicstyle=\ttfamily\scriptsize]
select name 
from instructor 
where salary > some ( 
          select salary 
          from instructor 
          where dept_name = 'Biology' 
      );
        \end{lstlisting}
    \end{itemize}
\end{frame}

\begin{frame}{Definition of 'some' Clause}
    \begin{itemize}
        \item $F \langle comp \rangle \textbf{ some } r \Leftrightarrow \exists t \in r \text{ such that } (F \langle comp \rangle t)$
        \vspace{0.5em}
        
        Where $\langle comp \rangle$ can be: $<, \le, >, \ge, =, \neq$
        \vspace{1em}
        \item \texttt{some} represents existential quantification
        \item Example: $\texttt{5 < some (0, 5, 6)} = \texttt{true}$
    \end{itemize}
\end{frame}

\begin{frame}[fragile]{Set Comparison - 'all' Clause}
    \begin{itemize}
        \item Find the names of all instructors whose salary is greater than the salary of \textbf{all} instructors in the Biology department
        \begin{lstlisting}[language=SQL]
select name 
from instructor 
where salary > all ( 
          select salary 
          from instructor 
          where dept_name = 'Biology' 
      );
        \end{lstlisting}
    \end{itemize}
\end{frame}

\begin{frame}{Definition of 'all' Clause}
    \begin{itemize}
        \item $F \langle comp \rangle \textbf{ all } r \Leftrightarrow \forall t \in r \text{ such that } (F \langle comp \rangle t)$
        \vspace{0.5em}
        
        Where $\langle comp \rangle$ can be: $<, \le, >, \ge, =, \neq$
        \vspace{1em}
        \item \texttt{all} represents universal quantification
        \item Example: $\texttt{5 < all (0, 5, 6)} = \texttt{false}$
        \item Example: $\texttt{5 < all (6, 10)} = \texttt{true}$
    \end{itemize}
\end{frame}

\begin{frame}{Test for Empty Relations: 'exists'}
    \begin{itemize}
        \item The \texttt{exists} construct returns the value \textbf{true} if the argument subquery is nonempty
        \vspace{1em}
        \item $\textbf{exists } r \Leftrightarrow r \neq \emptyset$
        \vspace{0.5em}
        \item $\textbf{not exists } r \Leftrightarrow r = \emptyset$
    \end{itemize}
\end{frame}

\begin{frame}[fragile]{Use of 'exists' Clause}
    \begin{itemize}
        \item Yet another way of specifying the query "Find all courses taught in both the Fall 2009 semester and in the Spring 2010 semester"
        \begin{lstlisting}[language=SQL, basicstyle=\ttfamily\scriptsize]
select course_id 
from section as S 
where semester = 'Fall' and year = 2009 and 
      exists ( 
          select * 
          from section as T 
          where semester = 'Spring' and year = 2010 
            and S.course_id = T.course_id 
      );
        \end{lstlisting}
        \item \textbf{Correlation name} - variable \texttt{S} in the outer query
        \item \textbf{Correlated subquery} - the inner query
    \end{itemize}
\end{frame}

\begin{frame}[fragile]{Use of 'not exists' Clause}
    \begin{itemize}
        \item Find all students who have taken all courses offered in the Biology department.
        \begin{lstlisting}[language=SQL, basicstyle=\ttfamily\scriptsize]
select distinct S.ID, S.name 
from student as S 
where not exists ( 
          (select course_id 
           from course 
           where dept_name = 'Biology') 
          except 
          (select T.course_id 
           from takes as T 
           where S.ID = T.ID) 
      );
        \end{lstlisting}
        \item First nested query lists all courses offered in Biology
        \item Second nested query lists all courses a particular student took
        \item Note: $X - Y = \emptyset \Leftrightarrow X \subseteq Y$
        \item Note: Cannot write this query using \texttt{= all} and its variants
    \end{itemize}
\end{frame}

\begin{frame}[fragile]{Test for Absence of Duplicate Tuples: 'unique'}
    \begin{itemize}
        \item The \texttt{unique} construct tests whether a subquery has any duplicate tuples in its result
        \item The \texttt{unique} construct evaluates to 'true' if a given subquery contains no duplicates
        \item Find all courses that were offered at most once in 2009
        \begin{lstlisting}[language=SQL, basicstyle=\ttfamily\scriptsize]
select T.course_id 
from course as T 
where unique ( 
          select R.course_id 
          from section as R 
          where T.course_id = R.course_id and R.year = 2009 
      );
        \end{lstlisting}
    \end{itemize}
\end{frame}

\section{Subqueries in the From Clause}

\begin{frame}
  \vfill
  \centering
  \begin{beamercolorbox}[sep=8pt,center,shadow=true,rounded=true]{title}
    \usebeamerfont{title}Subqueries in the From Clause\par%
  \end{beamercolorbox}
  \vfill
\end{frame}

\begin{frame}[fragile]{Subqueries in the From Clause}
    \begin{itemize}
        \item SQL allows a subquery expression to be used in the \texttt{from} clause
        \item Find the average instructors' salaries of those departments where the average salary is greater than \$42,000
        \begin{lstlisting}[language=SQL, basicstyle=\ttfamily\scriptsize]
select dept_name, avg_salary 
from (select dept_name, avg(salary) as avg_salary 
      from instructor 
      group by dept_name) 
where avg_salary > 42000;
        \end{lstlisting}
        \item Note that we do not need to use the \texttt{having} clause
        \item Another way to write above query:
        \begin{lstlisting}[language=SQL, basicstyle=\ttfamily\scriptsize]
select dept_name, avg_salary 
from (select dept_name, avg(salary) 
      from instructor 
      group by dept_name) 
      as dept_avg (dept_name, avg_salary) 
where avg_salary > 42000;
        \end{lstlisting}
    \end{itemize}
\end{frame}

\begin{frame}[fragile]{With Clause}
    \begin{itemize}
        \item The \texttt{with} clause provides a way of defining a temporary relation whose definition is available only to the query in which the \texttt{with} clause occurs
        \item Find all departments with the maximum budget
        \begin{lstlisting}[language=SQL]
with max_budget(value) as ( 
    select max(budget) 
    from department 
) 
select department.name 
from department, max_budget 
where department.budget = max_budget.value;
        \end{lstlisting}
    \end{itemize}
\end{frame}

\begin{frame}[fragile]{Complex Queries using With Clause}
    \begin{itemize}
        \item Find all departments where the total salary is greater than the average of the total salary at all departments
        \begin{lstlisting}[language=SQL, basicstyle=\ttfamily\scriptsize]
with dept_total (dept_name, value) as (
    select dept_name, sum(salary) 
    from instructor 
    group by dept_name
),
dept_total_avg(value) as (
    select avg(value) 
    from dept_total
) 
select dept_name 
from dept_total, dept_total_avg 
where dept_total.value > dept_total_avg.value;
        \end{lstlisting}
    \end{itemize}
\end{frame}

\section{Subqueries in the Select Clause}

\begin{frame}
  \vfill
  \centering
  \begin{beamercolorbox}[sep=8pt,center,shadow=true,rounded=true]{title}
    \usebeamerfont{title}Subqueries in the Select Clause\par%
  \end{beamercolorbox}
  \vfill
\end{frame}

\begin{frame}[fragile]{Scalar Subquery}
    \begin{itemize}
        \item Scalar subquery is one which is used where a single value is expected
        \item List all departments along with the number of instructors in each department
        \begin{lstlisting}[language=SQL]
select dept_name, 
       (select count(*) 
        from instructor 
        where department.dept_name = instructor.dept_name) 
        as num_instructors 
from department;
        \end{lstlisting}
        \item Runtime error if subquery returns more than one result tuple
    \end{itemize}
\end{frame}

\section{Modifications of the Database}

\begin{frame}
  \vfill
  \centering
  \begin{beamercolorbox}[sep=8pt,center,shadow=true,rounded=true]{title}
    \usebeamerfont{title}Modifications of the Database\par%
  \end{beamercolorbox}
  \vfill
\end{frame}

\begin{frame}{Modification of the Database}
    \begin{itemize}
        \item Deletion of tuples from a given relation
        \item Insertion of new tuples into a given relation
        \item Updating of values in some tuples in a given relation
    \end{itemize}
\end{frame}

\begin{frame}[fragile]{Deletion}
    \begin{itemize}
        \item Delete all instructors
        \begin{lstlisting}[language=SQL]
delete from instructor;
        \end{lstlisting}
        \item Delete all instructors from the Finance department
        \begin{lstlisting}[language=SQL]
delete from instructor where dept_name = 'Finance';
        \end{lstlisting}
        \item Delete all tuples in the \texttt{instructor} relation for those instructors associated with a department located in the Watson building
        \begin{lstlisting}[language=SQL]
delete from instructor 
where dept_name in (
          select dept_name 
          from department 
          where building = 'Watson'
      );
        \end{lstlisting}
    \end{itemize}
\end{frame}

\begin{frame}[fragile]{Deletion (2)}
    \begin{itemize}
        \item Delete all instructors whose salary is less than the average salary of instructors
        \begin{lstlisting}[language=SQL]
delete from instructor 
where salary < (
          select avg(salary) 
          from instructor
      );
        \end{lstlisting}
        \item \textbf{Problem}: as we delete tuples from \texttt{instructor}, the average salary changes
        \item \textbf{Solution used in SQL}:
        \begin{enumerate}
            \item First, compute \texttt{avg (salary)} and find all tuples to delete
            \item Next, delete all tuples found above (without recomputing avg or retesting the tuples)
        \end{enumerate}
    \end{itemize}
\end{frame}

\begin{frame}[fragile]{Insertion}
    \begin{itemize}
        \item Add a new tuple to \texttt{course}
        \begin{lstlisting}[language=SQL]
insert into course 
values ('CS-437', 'Database Systems', 'Comp. Sci.', 4);
        \end{lstlisting}
        \item or equivalently:
        \begin{lstlisting}[language=SQL]
insert into course (course_id, title, dept_name, credits) 
values ('CS-437', 'Database Systems', 'Comp. Sci.', 4);
        \end{lstlisting}
        \item Add a new tuple to \texttt{student} with \texttt{tot\_creds} set to null
        \begin{lstlisting}[language=SQL]
insert into student 
values ('3003', 'Green', 'Finance', null);
        \end{lstlisting}
    \end{itemize}
\end{frame}

\begin{frame}[fragile]{Insertion (2)}
    \begin{itemize}
        \item Add all instructors to the \texttt{student} relation with \texttt{tot\_creds} set to 0
        \begin{lstlisting}[language=SQL]
insert into student 
select ID, name, dept_name, 0 
from instructor;
        \end{lstlisting}
        \item The \texttt{select from where} statement is evaluated fully before any of its results are inserted into the relation
        \item Otherwise queries like
        \begin{lstlisting}[language=SQL]
insert into table1 
select * from table1;
        \end{lstlisting}
        would cause problem
    \end{itemize}
\end{frame}

\begin{frame}[fragile]{Updates}
    \begin{itemize}
        \item Increase salaries of instructors whose salary is over \$100,000 by 3\%, and all others by a 5\%
        \item Write two \texttt{update} statements:
        \begin{lstlisting}[language=SQL]
update instructor 
set salary = salary * 1.03 
where salary > 100000; 

update instructor 
set salary = salary * 1.05 
where salary <= 100000;
        \end{lstlisting}
        \item \textbf{The order is important}
        \item Can be done better using the \texttt{case} statement (next slide)
    \end{itemize}
\end{frame}

\begin{frame}[fragile]{Case Statement for Conditional Updates}
    \begin{itemize}
        \item Same query as before but with \texttt{case} statement
        \begin{lstlisting}[language=SQL]
update instructor 
set salary = case
    when salary <= 100000 then salary * 1.05 
    else salary * 1.03
end;
        \end{lstlisting}
    \end{itemize}
\end{frame}

\begin{frame}[fragile]{Updates with Scalar Subqueries}
    \begin{itemize}
        \item Recompute and update \texttt{tot\_creds} value for all students
        \begin{lstlisting}[language=SQL, basicstyle=\ttfamily\scriptsize]
update student S 
set tot_creds = (
    select sum(credits) 
    from takes, course 
    where takes.course_id = course.course_id 
      and S.ID = takes.ID 
      and takes.grade <> 'F' 
      and takes.grade is not null
);
        \end{lstlisting}
        \item Sets \texttt{tot\_creds} to null for students who have not taken any course
        \item Instead of \texttt{sum (credits)}, use:
        \begin{lstlisting}[language=SQL, basicstyle=\ttfamily\scriptsize]
case 
    when sum(credits) is not null then sum(credits) 
    else 0 
end
        \end{lstlisting}
    \end{itemize}
\end{frame}

\begin{frame}{Module Summary}
    \begin{itemize}
        \item Introduced nested subquery in SQL
        \item Introduced data modification
    \end{itemize}
    \vspace{2em}
    \begin{center}
    \small \textit{Slides used in this presentation are borrowed from \url{http://db-book.com/} with kind permission of the authors.} \\
    \small \textit{Edited and new slides are marked with 'PPD'.}
    \end{center}
\end{frame}

\end{document}
